\begin{problem}[Automorphisms of the affine line]\label{prob:vi19-p12}
Let $k$ be a field. Prove every $k$-algebra automorphism of $k[t]$ has $t\mapsto at+b$ with $a\in k^\times$, and translate this to automorphisms of $\A^1_k$ whose pullback fixes $k$.
\end{problem}

\ifdefined\FullProblemDossiers
\paragraph{Why this problem matters.}
This is a concrete classification using the affine anti-equivalence.

\paragraph{Useful ideas before starting.}
Use degrees of a polynomial and its inverse under composition.

\paragraph{Guided hints.}
\begin{enumerate}
\item Write $\varphi(t)=p(t)$ and $\varphi^{-1}(t)=q(t)$.
\item Use $q(p(t))=t$ to compare degrees.
\end{enumerate}

\begin{solution}
% VI19 full-solutions refinement v1
Let \(\varphi:k[t]\to k[t]\) be a \(k\)-algebra automorphism and set \(p(t)=\varphi(t)\).  Since \(\varphi\) is surjective, its inverse \(\psi\) is also a \(k\)-algebra automorphism; write \(q(t)=\psi(t)\).  Applying \(\psi\circ\varphi=\operatorname{id}\) to \(t\) gives
\[
q(p(t))=t.
\]
Neither \(p\) nor \(q\) is constant, because an automorphism cannot send \(t\) into \(k\).  Therefore
\[
1=\deg t=\deg(q\circ p)=(\deg q)(\deg p).
\]
Both degrees are positive integers, so \(\deg p=\deg q=1\).  Hence
\[
p(t)=at+b
\]
with \(a\in k^\times\) and \(b\in k\).

Conversely, every substitution \(t\mapsto at+b\) with \(a\neq0\) defines a \(k\)-algebra automorphism, whose inverse sends
\[
t\longmapsto a^{-1}(t-b).
\]
Under affine anti-equivalence, these ring automorphisms correspond exactly to automorphisms of \(\mathbb A_k^1=\Spec k[t]\) whose pullback fixes the base field \(k\).  On closed \(k\)-points the geometric map is the inverse affine-linear transformation dictated by contravariance; the ring classification is the intrinsic statement.
\end{solution}

\paragraph{What happened mathematically.}
Automorphisms of an affine scheme are ring automorphisms read in the opposite direction.

\paragraph{Extensions.}
Compare with automorphisms of $\A^n$ for $n>1$.

\paragraph{Provenance.}
\textbf{New canonical synthesis; no numbered legacy Problem is claimed.}
\else
\paragraph{Provenance.} \textbf{New canonical synthesis; no numbered legacy Problem is claimed.}
\fi
